\documentclass{article}
\usepackage{amsmath}
\usepackage{amssymb}
\usepackage{amsthm}
\usepackage{amsfonts}
\usepackage{amscd}
\usepackage{amsthm}
\usepackage{amsbsy}
\usepackage{amsthm}
\theoremstyle{definition}
\newtheorem{theorem}{Theorem}
\newtheorem{lemma}{Lemma}
\newtheorem{proposition}{Proposition}
\newtheorem{corollary}{Corollary}
\newtheorem{definition}{Definition}
\newtheorem{example}{Example}
\newtheorem{remark}{Remark}
\newtheorem{assumption}{Assumption}

\begin{document}

\newcommand{\res}{\mathbin{\backslash}}

\begin{theorem}
For every \(x\),
\[
((((x\res e)\res e)\res x)\res x)
=
((((((x\res e)\res e)\res x)\res x)\res e)\res e).
\]
\end{theorem}

\begin{proof}
Introduce the abbreviations
\[
u=x\res e,\qquad
v=u\res e,\qquad
w=v\res x,
\]
and
\[
a=w\res x,\qquad
b=a\res e,\qquad
c=b\res e.
\]
It is therefore enough to prove \(a=c\).

First, the contraction and mingle axioms give
\[
p\res(p\res q)\leq p\res q
\quad\text{and}\quad
p\res q\leq p\res(p\res q).
\]
Hence, by antisymmetry,
\begin{equation}
p\res(p\res q)=p\res q.
\tag{1}
\end{equation}

We shall also use the following consequence of the axioms:
\begin{equation}
p\leq q
\quad\Longrightarrow\quad
q\res e\leq p\res e.
\tag{2}
\end{equation}
Indeed, from
\[
q\leq(q\res e)\res e
\]
and \(p\leq q\), transitivity gives
\[
p\leq(q\res e)\res e.
\]
Using \(r\leq s\iff e\leq r\res s\), followed by exchange, we obtain
\[
e\leq p\res\bigl((q\res e)\res e\bigr)
  =(q\res e)\res(p\res e).
\]
Applying the same order characterization once more proves (2).

The axiom \(p\leq(p\res q)\res q\), with \(p=a\) and \(q=e\),
gives
\begin{equation}
a\leq(a\res e)\res e=c.
\tag{3}
\end{equation}

We next prove
\begin{equation}
a=w\res v.
\tag{4}
\end{equation}
The axiom
\[
p\res q\leq(q\res r)\res(p\res r)
\]
with \(p=w,q=x,r=e\) gives
\[
a=w\res x
 \leq(x\res e)\res(w\res e)
 =u\res(w\res e).
\]
By exchange,
\[
u\res(w\res e)=w\res(u\res e)=w\res v,
\]
so \(a\leq w\res v\).

Conversely, applying the same axiom with \(p=w,q=v,r=x\) gives
\[
w\res v
 \leq(v\res x)\res(w\res x)
 =w\res a.
\]
Since \(a=w\res x\), equation (1) yields
\[
w\res a=w\res(w\res x)=w\res x=a.
\]
Thus \(w\res v\leq a\), proving (4) by antisymmetry.

From (3) and (2),
\[
c\res e\leq a\res e=b.
\]
On the other hand,
\[
b\leq(b\res e)\res e=c\res e.
\]
Consequently,
\begin{equation}
b=c\res e.
\tag{5}
\end{equation}

Using exchange, first with \(b=a\res e\) and then with
\(b=c\res e\), we obtain
\[
u\res b
 =u\res(a\res e)
 =a\res(u\res e)
 =a\res v
\]
and
\[
u\res b
 =u\res(c\res e)
 =c\res(u\res e)
 =c\res v.
\]
Therefore
\begin{equation}
a\res v=c\res v.
\tag{6}
\end{equation}

Since \(a=w\res v\), exchange gives, for every \(t\),
\[
t\res a
 =t\res(w\res v)
 =w\res(t\res v).
\]
Taking \(t=a\) and \(t=c\), and using (6), gives
\[
a\res a
 =w\res(a\res v)
 =w\res(c\res v)
 =c\res a.
\]
The axiom \(e\leq a\res a\) therefore implies
\[
e\leq c\res a.
\]
By \(p\leq q\iff e\leq p\res q\), this means
\[
c\leq a.
\]
Together with (3), antisymmetry yields \(a=c\), which is precisely
the required identity.
\end{proof}

\end{document}